% ============================================================================
% CAPÍTULO 4 — JavaScript\index{JavaScript}: a linguagem da web
% ============================================================================
\chapter{JavaScript: A Linguagem da Web}
\label{cap:js}

\begin{caixaconceito}
\textbf{JavaScript} é a linguagem de \emph{comportamento} da web — a única que
os navegadores executam nativamente. Hoje ela também roda em servidores
(Node.js), em bancos e até em edge computing: é a linguagem mais usada do
mundo \cite{so2025}.
\end{caixaconceito}

Ao final deste capítulo, você será capaz de:
\begin{objetivos}
  \item Escrever código JavaScript moderno (ES2023+) com \texttt{const}, \texttt{let}, arrow functions e template literals;
  \item Manipular arrays e objetos com \texttt{map}, \texttt{filter} e \texttt{reduce};
  \item Interagir com a página via DOM\index{DOM} e eventos;
  \item Programar de forma assíncrona com \texttt{async/await} e \texttt{fetch};
  \item Aplicar algoritmos úteis de mercado (embaralhamento Fisher--Yates, cronômetros, recordes em \texttt{localStorage});
  \item Entregar o projeto do capítulo: um jogo da memória completo.
\end{objetivos}

\section{O papel do JavaScript}

HTML diz \emph{o que é}, CSS diz \emph{como parece}, JavaScript diz \emph{o que
faz}: valida formulários, busca dados em APIs, anima interfaces e atualiza a
página sem recarregá-la. É o motor de interatividade da web moderna — e a base
de tudo o que você construirá nos capítulos seguintes (React, Node.js, Next.js).

\begin{caixadica}
JavaScript \textbf{não} é Java. A confusão é histórica (marketing dos anos 90),
mas as linguagens não têm relação. JavaScript é dinâmica, tipada
\emph{implicitamente} e executada por engines como o V8 (Chrome/Node). No
capítulo~\ref{cap:typescript} você adicionará tipos estáticos com TypeScript.
\end{caixadica}

\section{Variáveis e tipos primitivos}

\begin{lstlisting}[style=livro, language=JavaScript,
  caption={Declarações e tipos primitivos.}, label={list:js-variaveis}]
const nome = "Ana";          // constante — não pode ser reatribuída
let idade = 28;              // variável mutável
var legado = "evite";        // var tem escopo de função; quase nunca use

const ativo = true;          // boolean
let saldo = 19.9;            // number (inteiros e decimais são o mesmo tipo)
const nulo = null;           // ausência intencional de valor
let indefinido;              // undefined: declarada sem valor
const simbolo = Symbol("id");// identificador único
const big = 9007199254740993n; // BigInt para inteiros gigantes

console.log(typeof nome);    // "string"
\end{lstlisting}

\begin{caixaarmadilha}
Use \texttt{const} por padrão e \texttt{let} apenas quando precisar
reatribuir. Isso torna o código previsível e evita bugs de estado. O \texttt{var}
fica para código legado (você o encontrará em sistemas antigos no mercado).
\end{caixaarmadilha}

\section{Igualdade, template literals e estruturas de controle}

\begin{lstlisting}[style=livro, language=JavaScript,
  caption={Igualdade estrita, template literals e controle de fluxo.},
  label={list:js-controle}]
// === compara valor E tipo (sempre use ===, nunca ==)
console.log(5 === "5");   // false
console.log(5 == "5");    // true (coerção implícita — evite!)

// Template literals com crases
const cidade = "Belo Horizonte";
console.log(`Bem-vindo, ${nome}, de ${cidade}!`);

// Controle de fluxo
if (idade >= 18) {
  console.log("Maior de idade");
} else {
  console.log("Menor de idade");
}

for (let i = 1; i <= 3; i++) {
  console.log(i);
}

const cores = ["azul", "verde"];
for (const cor of cores) {
  console.log(cor);
}
\end{lstlisting}

\section{Funções: declarações, expressões e arrow functions}

\begin{lstlisting}[style=livro, language=JavaScript,
  caption={Os três formatos de função.}, label={list:js-funcoes}]
// 1. Declaração (hoisted — pode ser chamada antes da definição)
function calcularImposto(valor, aliquota) {
  return valor * aliquota;
}

// 2. Expressão de função
const calcularJuros = function (valor, taxa) {
  return valor * (1 + taxa);
};

// 3. Arrow function (padrão moderno)
const dobrar = (x) => x * 2;
const somar = (a, b) => {
  const total = a + b;
  return total;
};

console.log(calcularImposto(1000, 0.15)); // 150
\end{lstlisting}

\begin{caixaconceito}
\textbf{Escopo} define onde uma variável é visível. \texttt{let} e
\texttt{const} têm escopo de bloco (\texttt{\{...\}}); \texttt{var} tem escopo
de função. Arrow functions não criam seu próprio \texttt{this} — o que evita
uma das maiores dores de cabeça do JavaScript antigo.
\end{caixaconceito}

\section{Arrays e objetos: o trio \texttt{map}, \texttt{filter}, \texttt{reduce}}

Estas três funções resolvem a maioria dos problemas de processamento de dados
— você as usará todos os dias na carreira:

\begin{lstlisting}[style=livro, language=JavaScript,
  caption={map, filter e reduce — as três funções essenciais.},
  label={list:js-map}]
const pedidos = [
  { id: 1, produto: "Notebook", valor: 4500, pago: true },
  { id: 2, produto: "Mouse", valor: 120, pago: false },
  { id: 3, produto: "Teclado", valor: 350, pago: true },
];

// map: transforma cada item
const nomes = pedidos.map((p) => p.produto);
// ["Notebook", "Mouse", "Teclado"]

// filter: seleciona itens
const pagos = pedidos.filter((p) => p.pago);
// [Notebook, Teclado]

// reduce: acumula em um único valor
const totalPago = pedidos
  .filter((p) => p.pago)
  .reduce((soma, p) => soma + p.valor, 0);
// 4850

// Destructuring e spread
const { produto, valor } = pedidos[0];
const copia = [...pedidos, { id: 4, produto: "Monitor", valor: 900, pago: false }];
\end{lstlisting}

\begin{caixadica}
Pense em \texttt{map} como ``transforme cada item'', em \texttt{filter} como
``mantenha apenas os que passam no teste'' e em \texttt{reduce} como ``acumule
tudo em um resultado''. Esses três verbos substituem 90\% dos loops
\texttt{for} que você veria em código antigo.
\end{caixadica}

\section{O DOM: o JavaScript conversa com a página}

O \textbf{DOM} (Document Object Model) é a representação da página em árvore de
objetos que o JavaScript pode ler e alterar:

\begin{lstlisting}[style=livro, language=JavaScript,
  caption={Manipulação básica do DOM.}, label={list:js-dom}]
// Selecionar
const titulo = document.querySelector("h1");
const botoes = document.querySelectorAll(".botao");
const formulario = document.getElementById("cadastro");

// Ler e alterar
titulo.textContent = "Novo título";
titulo.style.color = "#1f4e79";
titulo.classList.add("destaque");

// Criar e inserir
const item = document.createElement("li");
item.textContent = "Item novo";
document.querySelector("ul").appendChild(item);

// Remover
titulo.remove();
\end{lstlisting}

\section{Eventos: a página reage ao usuário}

\begin{lstlisting}[style=livro, language=JavaScript,
  caption={Eventos: clique, input e envio de formulário.}, label={list:js-eventos}]
const botao = document.querySelector("#enviar");
const campo = document.querySelector("#nome");

botao.addEventListener("click", (evento) => {
  console.log("Clicado!", evento.target);
});

campo.addEventListener("input", () => {
  campo.dataset.tamanho = String(campo.value.length);
});

formulario.addEventListener("submit", (evento) => {
  evento.preventDefault(); // impede o recarregamento da página
  console.log("Enviando dados...");
});
\end{lstlisting}

\begin{caixaarmadilha}
Em um formulário, se o handler de \texttt{submit} não chamar
\texttt{evento.preventDefault()}, a página recarrega e você perde o estado —
o bug mais comum de iniciantes que ``tentam fazer fetch\index{fetch}''. Esse detalhe
aparece em toda entrevista de frontend.
\end{caixaarmadilha}

\section{Assincronismo: promises, \texttt{async/await} e \texttt{fetch}}

O JavaScript é \emph{mono-thread} com um \emph{event loop}: ele não espera
operações lentas (rede, disco) paradas; ele agenda uma continuação e segue
executando. As \textbf{promises} modelam esse futuro, e
\texttt{async/await} o torna legível:

\begin{lstlisting}[style=livro, language=JavaScript,
  caption={fetch + async/await: consumindo uma API pública.},
  label={list:js-fetch}]
async function carregarUsuarios() {
  try {
    const resposta = await fetch("https://api.exemplo.com/usuarios");
    if (!resposta.ok) {
      throw new Error(`Erro HTTP ${resposta.status}`);
    }
    const usuarios = await resposta.json();
    return usuarios;
  } catch (erro) {
    console.error("Falha ao carregar:", erro);
    return [];
  }
}

// Uso
const usuarios = await carregarUsuarios();
console.log(usuarios.length);
\end{lstlisting}

\begin{caixaconceito}
\texttt{await} só funciona dentro de funções \texttt{async}. Toda função
\texttt{async} retorna uma \textbf{Promise}. O fluxo \emph{try/catch/finally}
é o padrão para tratar erros de rede — você o usará em todas as APIs que
construir a partir do capítulo~\ref{cap:http}.
\end{caixaconceito}

\section{Módulos ES: organizando o código}

\begin{lstlisting}[style=livro, language=JavaScript,
  caption={Módulos ES: exportar e importar.}, label={list:js-modulos}]
// utils.js
export function formatarMoeda(valor) {
  return valor.toLocaleString("pt-BR", { style: "currency", currency: "BRL" });
}

export const TAXA_PADRAO = 0.15;

// main.js
import { formatarMoeda, TAXA_PADRAO } from "./utils.js";
console.log(formatarMoeda(1234.5)); // R$ 1.234,50
\end{lstlisting}

\begin{caixadica}
Módulos ES (\texttt{import}/\texttt{export}) são o padrão moderno em todos os
frameworks. A partir do capítulo~\ref{cap:node}, você os usará diariamente —
inclusive para testar funções isoladas com Vitest.
\end{caixadica}

% ============================================================================
\section{Projeto do capítulo: jogo da memória}
\label{sec:projeto4}

\begin{caixaprojeto}
\textbf{Objetivo:} construir um jogo da memória completo em JavaScript puro,
com embaralhamento aleatório, cronômetro, contador de jogadas e recorde
persistente no navegador.

\textbf{Critérios de aceite:}
\begin{itemize}[leftmargin=*, itemsep=0.2em]
  \item Grade 4×4 com 8 pares de emojis embaralhados a cada partida;
  \item Algoritmo de embaralhamento \textbf{Fisher--Yates} implementado do zero;
  \item Lógica de pares: cartas erradas voltam após 700ms; certas ficam verdes;
  \item Cronômetro e contador de jogadas atualizados em tempo real;
  \item Recorde salvo em \texttt{localStorage} e exibido na tela;
  \item Acessível: botões reais, \texttt{aria-label} dinâmico e \texttt{:focus-visible}.
\end{itemize}
\end{caixaprojeto}

O código completo do jogo (Listagem~\ref{list:jogo-js}):

\lstinputlisting[style=livro, language=JavaScript,
  caption={Lógica completa do jogo — \texttt{codigo/cap04-jogo-memoria/script.js}.},
  label={list:jogo-js}]
  {\codigopath/cap04-jogo-memoria/script.js}

\begin{caixadica}
Para rodar: abra \texttt{codigo/cap04-jogo-memoria/index.html} direto no
navegador, ou sirva com \texttt{npx serve codigo/cap04-jogo-memoria}.
Teste também no DevTools $\rightarrow$ Console: os \texttt{console.log} ajudam
a depurar passo a passo. \emph{Nota de tipografia}: os emojis das cartas
podem aparecer como símbolos simples no PDF conforme a fonte instalada;
no navegador eles aparecem coloridos, como o jogo foi desenhado.
\end{caixadica}

\section{Exercícios de fixação}

\begin{exercicios}
  \item Qual a diferença entre \texttt{const}, \texttt{let} e \texttt{var}?
  Escreva um exemplo em que \texttt{let} é necessário.
  \item Por que usar \texttt{===} em vez de \texttt{==}? Dê um exemplo em que
  \texttt{==} produz um resultado surpreendente.
  \item Dado \texttt{const nums = [1, 2, 3, 4, 5]}, use \texttt{map} para
  dobrar, \texttt{filter} para pares e \texttt{reduce} para somar — sem loops.
  \item O que faz \texttt{evento.preventDefault()} em um formulário?
  \item Explique o que é uma Promise\index{Promise} e por que \texttt{async/await} foi criado.
  \item Crie um botão que, ao ser clicado, adiciona um item a uma lista
  \texttt{ul} usando \texttt{createElement} e \texttt{appendChild}.
\end{exercicios}

\section{Desafios}

\begin{desafios}
  \item \textbf{FizzBuzz}: escreva FizzBuzz (múltiplos de 3 $\rightarrow$ Fizz,
  5 $\rightarrow$ Buzz, ambos $\rightarrow$ FizzBuzz) com arrow function e \texttt{map}.
  \item \textbf{Contador de palavras}: dado um texto, retorne um objeto com a
  frequência de cada palavra usando \texttt{reduce}.
  \item \textbf{Cronômetro Pomodoro}: adapte o cronômetro do jogo em um app de
  25 minutos com alerta sonoro (Web Audio API) e pausa/retomada.
  \item \textbf{Modo difícil}: aumente o jogo para 6×6 (18 pares) com
  dificuldade selecionável e recorde por dificuldade.
\end{desafios}

\section{Exercícios de nível sênior}

\begin{seniores}
  \item \textbf{Debounce:} implemente \texttt{debounce(fn, 300)} — função que
  só executa 300ms após a última chamada — e use-a em uma busca que faz
  \texttt{fetch} a cada tecla. Explique por que isso importa para performance.
  \item \textbf{Event delegation:} refatore a lista de cartas do jogo para usar
  \emph{um único} listener na grade (\texttt{event.target}), em vez de um
  listener por carta. Compare a performance com 1000 cartas.
  \item \textbf{Closures em produção:} crie um contador usando closure (função
  que retorna função) e explique onde closures\index{closures} aparecem em React (pistas:
  \texttt{useState}, \texttt{useEffect}).
  \item \textbf{Consumo real de API:} busque dados da API pública
  \url{https://api.github.com/users/usuario} e renderize nome, avatar e
  repositórios de um usuário, com tratamento de erro para \texttt{404}.
\end{seniores}

\section{Armadilhas comuns}

\begin{itemize}[leftmargin=*, itemsep=0.4em]
  \item Esquecer \texttt{await} (você recebe uma Promise, não o dado);
  \item Usar \texttt{==} e sofrer com coerção implícita;
  \item Mutar arrays/objetos por engano (prefira \texttt{map}/\texttt{filter}/\texttt{spread});
  \item Anexar 1000 listeners onde 1 (delegation) resolve;
  \item Bloquear a UI com loops síncronos pesados — use \texttt{setTimeout}/\texttt{requestAnimationFrame}.
\end{itemize}

\section{Resumo do capítulo}

\begin{itemize}[leftmargin=*, itemsep=0.3em]
  \item JavaScript é a linguagem de comportamento da web e a mais usada do mundo;
  \item \texttt{const}/\texttt{let}, \texttt{===}, template literals e arrow functions são o padrão moderno;
  \item \texttt{map}, \texttt{filter}, \texttt{reduce} resolvem o processamento de dados;
  \item DOM + eventos conectam o código à interface;
  \item \texttt{async/await} + \texttt{fetch} conversam com APIs;
  \item Projeto entregue: jogo da memória com Fisher--Yates, cronômetro e recorde.
\end{itemize}

\section{Referências oficiais para aprofundar}

\begin{itemize}[leftmargin=*, itemsep=0.3em]
  \item MDN — JavaScript: \url{https://developer.mozilla.org/pt-BR/docs/Web/JavaScript}
  \item MDN — Guia de JavaScript: \url{https://developer.mozilla.org/en-US/docs/Web/JavaScript/Guide}
  \item JavaScript.info (referência didática): \url{https://javascript.info}
  \item MDN — Fetch API: \url{https://developer.mozilla.org/en-US/docs/Web/API/Fetch_API}
  \item ECMAScript Language Specification: \url{https://tc39.es/ecma262/}
\end{itemize}
